% =====================================================================================================
\section{The assumptions}\label{sec:assumptions}\label{sec:standing}

This section states every assumption of the main theorems once, and judges it. Table~\ref{tab:a-summary} gives the
verdicts at a glance. Then each assumption has one entry with five parts: \textbf{(a)} the statement; \textbf{(b)} what it
means in plain words; \textbf{(c)} what it rules out; \textbf{(d)} whether it is reasonable, for real designs, with an
example where one exists; and \textbf{(e)} how it compares with \citet{SongLiang2015}: the same, weaker (asks less),
stronger (asks more), or without an analogue. Page numbers for \SL{} are journal pages; ``supp.\ p.'' is a page of their
online supplement; $\bar r_n$ is their cap on model size.

\paragraph{Our designs.} Part (d) refers to the designs of our simulation studies. All have i.i.d.\ Gaussian rows, noise
variance $1$ and fixed nonzero coefficients, in three classes: \emph{independent} (S1 of the 2,400-fit study; IND of the
$10^6$ plan); \emph{locally correlated} (S2, AR(1) at $0.5$; S3 and S4, blocks of five at $0.7$ and $0.85$; S6, blocks of
five at $0.4$--$0.9$; AR, AR(1) at $0.6$; BLK, blocks of $20$ at $0.5$); and \emph{pervasive factor} (S5, blocks of five
at $0.7$ plus a global factor at $0.3$; the strong-factor study, one factor at $0.3$, $0.5$ or $0.7$ or two factors, with
all true effects of one sign, $n$ from $50$ to $150$ and $p$ up to $10n$; FAC, one factor at $0.3$). The 2,400-fit study has
$n=200$, $p\in\{500,1000\}$ and $14$ true predictors, one in each group of five; the $10^6$ plan has $5$ true predictors.
A design is \emph{inside} an assumption when the design, continued as a sequence in $n$ with its structure and the number
of true predictors fixed, satisfies it; we judge the qualitative conditions, not the conservative constants of the proofs.

\paragraph{The short answer.} The posterior theorems (Results 1 and 2) use the same kinds of assumptions as \SL{}, and
they hold on all our designs, strong common factors included. The theorem about the reported model (Result 3) adds two
conditions that \SL{} do not have, visibility and a small leakage set; they hold for block and autoregressive designs and
fail for pervasive factors. That is why Result 3 is stated for block and AR designs, and why Result 4 exists: it needs
neither condition.

\begin{table}[!t]
\centering\footnotesize
\setlength{\tabcolsep}{3pt}
\renewcommand{\arraystretch}{1.05}
\caption{The assumptions at a glance. ``Used by'' names the results (Section~\ref{sec:proved}). The last column compares
with \citet{SongLiang2015}: \emph{same}, \emph{weaker} (ours asks less), \emph{stronger} (ours asks more), \emph{mixed}
(depends on the regime), or \emph{none} (no analogue).}\label{tab:a-summary}
\begin{tabular}{@{}>{\raggedright\arraybackslash}p{0.2\linewidth}>{\raggedright\arraybackslash}p{0.08\linewidth}>{\raggedright\arraybackslash}p{0.38\linewidth}>{\raggedright\arraybackslash}p{0.28\linewidth}@{}}
\toprule
Assumption & Used by & Reasonable? & vs \SL \\
\midrule
\multicolumn{4}{@{}l}{\textbf{The data model}} \\
\lab{GL} Gaussian linear model & 1, 3, 4, 6 & the standard model; true by construction in our designs & same: (1), p.~947 \\
\lab{DIM} $p\le n^A$ & 1--3 & any polynomial $p$; $p=10^6$ at $n=200$ is $A\approx2.6$ & mixed: narrower at the top, wider at the bottom \\
\lab{SP} fixed true model size & 1, 3 & few strong effects; nearly necessary for the rate & same as (4), p.~950 \\
\midrule
\multicolumn{4}{@{}l}{\textbf{The predictors and the signal}} \\
\lab{SGR}, \lab{GR} law of the rows & 1--3 & Gaussian or sub-Gaussian; bounded features included, rare binary features not & stronger than Cond.~1, p.~953 (standardised predictors) \\
\lab{FSb}, \lab{FSt} fixed signals & 1, 3 & effects that do not shrink with $n$; Result 2 needs none & same, if their $\bstar$ is fixed \\
\lab{EFS} eigenvalue floor & 1, 3 & holds in all our designs (floors $0.1$ to $1$); no upper bound, so factors allowed & mixed: stronger unless their $\bar r_n$ is large \\
\ref{asm:dist}(d) no near-collinearity & 2 & only exact duplicates and exact dependences fail & mixed: none in their Thms 1--2 \\
\midrule
\multicolumn{4}{@{}l}{\textbf{The reported model (Result 3 only)}} \\
\lab{V} visibility & 3 & holds with aligned signs or generic coefficients; can fail by cancellation & stronger than Cond.~5, p.~958 \\
\ref{asm:reported}(a) small leakage set & 3 & holds for block and AR designs; fails for pervasive factors & none \\
\ref{asm:reported}(b),(c) floor and local regressions & 3 & follow from a floor on $\Sig$ and local correlation & floor as \lab{EFS}; locality none \\
\midrule
\multicolumn{4}{@{}l}{\textbf{Certified screening (Result 4 only)}} \\
floor on $k_n+\sz$ sets, signal bounds & 4 & fixed design; the sample-size budget is conservative & as \lab{EFS} and signals \\
\midrule
\multicolumn{4}{@{}l}{\textbf{Not statistical}} \\
conventions \as{C1}, \as{C2p}, \as{C3} & 3--5 & idealise the computer & presupposed by \SL{} too \\
\bottomrule
\end{tabular}
\end{table}

\newcommand{\Astate}{\par\smallskip\noindent\textbf{(a) Statement.}\ }
\newcommand{\Awords}[1]{\par\noindent\textbf{(b) In words.} #1}
\newcommand{\Arules}[1]{\par\noindent\textbf{(c) Rules out.} #1}
\newcommand{\Areas}[1]{\par\noindent\textbf{(d) Reasonable?} #1}
\newcommand{\Asl}[1]{\par\noindent\textbf{(e) vs \SL.} #1\par\medskip}

% ------------------------------------------------------------------------------------------------
\subsection{The data model}\label{sec:a-all}

\ahead{Gaussian linear model}{\lab{GL}}{}
\Astate
\begin{assumption}{GL}{Gaussian linear model}
The data follow~\eqref{eq:model}, with a fixed noise variance $\sigs^2>0$ and noise independent of the design.
\end{assumption}
\Awords{The response is a linear function of the predictors plus independent Gaussian noise of constant variance.}
\Arules{Heavy-tailed or heteroscedastic errors, and a nonlinear mean.}
\Areas{It is the working model of almost all theory for Bayesian variable selection, and it holds by construction in our
simulation designs. For real data it is a modelling choice, as in any linear-model analysis; Result~2 lets the noise level
vary with $n$.}
\Asl{\emph{The same} (their model (1), p.~947, and density (2), p.~950); ours adds an intercept.}

\ahead{Dimension}{\lab{DIM}}{}
\Astate
\begin{assumption}{DIM}{Dimension}
$2\le p\le n^A$ for a fixed $A\ge1$.
\end{assumption}
\Awords{The number of predictors may be much larger than $n$, but at most a fixed power of it.}
\Arules{Numbers of predictors that grow faster than every power of $n$, such as $\log p\asymp n^\alpha$.}
\Areas{Yes for the problems we target: $p=10^4$ to $10^6$ with $n$ in the hundreds. The largest exponent in our designs is
the $10^6$ plan at $n=200$, $A=\log10^6/\log200\approx2.6$; $p=10^7$ at $n=200$ would be $A\approx3.0$. Since $A$ may be
any fixed number, the assumption only excludes super-polynomial growth.}
\Asl{\emph{Mixed.} At the top ours is \emph{stronger}: their Corollary~1 allows $P_n<\exp(Cn^\alpha)$ (p.~954), and their
Theorem~S allows slightly more than polynomial growth (Corollary~S$_1$, $P_n<n^{\kappa\log n}$, supp.\ p.~16). At the
bottom ours is \emph{weaker}: it allows bounded $p$, while they need $P_n\succ n^\theta$ (p.~950).}

\ahead{Fixed true model size}{\lab{SP}}{}
\Astate
\begin{assumption}{SP}{Sparsity}
The true model size $\sz\ge1$ is fixed.
\end{assumption}
\Awords{A fixed, small number of predictors truly matter, however large $n$ and $p$ become.}
\Arules{True models whose size grows with $n$.}
\Areas{Yes when a few strong effects drive the response; it is the usual assumption for exact selection. It is also nearly
necessary for the rate we prove: on every design, $\pis{\Sz}\le\{1+(\sz+1)/(3\sqrt n)\}^{-1}$ (Result~6(a)), so a posterior
miss of order $\rn$ is possible only if $\sz=O(\sqrt{\log n})$. Result~4 does not assume it.}
\Asl{\emph{The same} as their strict sparsity (4), p.~950, which their Theorems~2, 3 and~S use; their Theorem~1 also
allows growing supports (p.~950), so against that theorem ours is \emph{stronger}.}

% ------------------------------------------------------------------------------------------------
\subsection{The predictors and the signal}\label{sec:a-post}

\ahead{Law of the rows}{\lab{SGR}, \lab{GR}}{}
\Astate
\begin{assumption}{GR}{Gaussian rows}
The rows of $\Xr$ are i.i.d.\ $\Normal_p(m_n,\Sig)$, where $\Sig$ is a correlation matrix.
\end{assumption}
\begin{assumption}{SGR}{Sub-Gaussian rows on sparse directions}
The raw rows are i.i.d.\ copies of a random vector $x\in\R^{\praw}$ with mean $m_n$ and covariance matrix $\Sig$, a
correlation matrix, and, for a fixed $L>0$, $\E\exp\{\lambda u^\top(x-m_n)\}\le\exp\{L^2\lambda^2u^\top\Sig u/2\}$ for every
$\lambda\in\R$ and every $u\in\R^{\praw}$ with at most $5(A+1)(\sz+1)\log n$ nonzero entries. Gaussian rows \as{GR} are
the case $L=1$.
\end{assumption}
\Awords{The observations are independent, and every sparse combination of predictors has Gaussian-like tails. Result~3
asks the rows to be exactly Gaussian; Results~1 and~2 ask only light tails.}
\Arules{Heavy-tailed predictors, and binary features whose success probability tends to zero (Results 1 and 2);
non-Gaussian rows (Result 3).}
\Areas{For Results~1 and~2, yes: \as{SGR} holds for Gaussian rows and for correlated mixtures $x=m_n+B\xi$ ($BB^\top$ a
correlation matrix) of independent features confined to $K$ standard deviations of their mean, with one common $K$; so
bounded data such as genotypes with allele frequencies bounded away from zero are inside. For Result~3 the Gaussian
requirement is a real restriction: such genotype data are outside Result~3 as stated, though the fixed-design version
(supplement, Theorem~\ref{thm:reported-fixed}) can cover an observed design that meets its conditions.}
\Asl{\emph{Stronger} than their Condition~1 (p.~953) on centred, unit-variance predictors: ours implies theirs.
Against their bounded predictors $\abs{x_i}\le1$ (Theorem~1, p.~952), neither implies the other.}

\ahead{Fixed signals}{\lab{FSb}, \lab{FSt}}{}
\Astate
\begin{assumption}{FSb}{Fixed signals, raw scale}
$b_{\mathrm{lo}}\le\abs{\beta^*_j}\le b_{\mathrm{hi}}$ for every $j\in\Sz$, with fixed $0<b_{\mathrm{lo}}\le b_{\mathrm{hi}}<\infty$.
\end{assumption}
\begin{assumption}{FSt}{Fixed signals, standardised scale}
$\yr=\alpha\one+X\tstar+\varepsilon$ with $\varepsilon\sim\Normal_n(0,\sigs^2I_n)$, a fixed $\sigs^2>0$, noise independent
of $\Xr$, a deterministic $\tstar$, $\sz\ge1$ fixed, and $\theta_{\mathrm{lo}}\le\abs{\theta^*_j}\le\theta_{\mathrm{hi}}$ for
$j\in\Sz$, with fixed $0<\theta_{\mathrm{lo}}\le\theta_{\mathrm{hi}}$.
\end{assumption}
\Awords{Each true effect is bounded away from zero and from infinity, by constants that do not change with $n$. \as{FSt}
fixes the effects on the standardised scale the method works on; it replaces \as{GL} and \as{SP} in Result~3.}
\Arules{True effects that shrink with $n$ (for example of order $\rn$), and effects that grow without bound.}
\Areas{Yes for selection: a predictor whose effect shrinks to zero cannot be selected reliably by any method. Result~2
needs no lower bound at all, so prediction is protected when some effects are weak; the fixed-design version of Result~1
(supplement, Theorem~\ref{thm:post-fixed}) lets the effects shrink down to the usual order $\sigs\rn$.}
\Asl{\emph{The same} if their $\bstar$ does not change with $n$, the natural reading of p.~950: they state no lower bound,
and fixed nonzero coefficients are what their theorems then ask. If their $\bstar$ may vary with $n$, their proofs need
rates (supp.\ pp.~6, 15), and our fixed signals ask more than those steps use.}

\ahead{Eigenvalue floor on small sets}{\lab{EFS}}{}
\Astate
\begin{assumption}{EFS}{Population eigenvalue floor}
$\lmin(\Sigma_{n,S})\ge\varphi^2$ for every $S$ with $\abs{S}\le5(A+1)(\sz+1)\log n$, for a fixed $\varphi\in(0,1]$ (for
instance, $\lmin(\Sig)\ge\varphi^2$).
\end{assumption}
\Awords{No small group of predictors, of up to logarithmic size, is close to linearly dependent in the population.}
\Arules{Exact or asymptotic duplicates of a predictor, and correlations that tend to one.}
\Areas{Yes, and it holds in every design we study. No upper bound on the eigenvalues is asked, so strong common factors
are allowed. The smallest eigenvalue of $\Sig$ is $1$ (S1, IND), $1/3$ (S2), $0.3$ (S3, S5), $0.15$ (S4), $0.1$ (S6),
$1-\rho$ (one factor at $\rho$), $0.40$ (two factors), $0.25$ (AR), $0.5$ (BLK) and $0.7$ (FAC). In real data with
near-duplicate predictors (two genes correlated at $0.99$) the floor is small and the constants degrade, but the
assumption still holds if the floor stays fixed.}
\Asl{\emph{Mixed.} Where their sets are no larger than ours, our fixed floor implies their (22)--(23) (supp.\ p.~13) and their
Condition~3 (p.~955), so ours is \emph{stronger}; their Condition~3 also asks only a vanishing level. Where their sets are
larger ($\bar r_n\succ\log n$), the two constrain different sets and neither implies the other.}

% ------------------------------------------------------------------------------------------------
\subsection{The assumptions of the distribution result}\label{sec:a-dist}

Result~2 has its own, weaker, setting. Its observations are $n$ i.i.d.\ copies $(x^{(i)},y^{(i)})$ of a pair $(x,y)$, with
$x\in\R^{\praw}$ the vector of all observed predictors, and $\zeta_j:=(x_j-\E x_j)/\sd(x_j)$. Fix a \emph{level} $M\ge1$;
the \emph{window} $q_n:=\min\{\praw,\lceil C_qM\log n\rceil\}$, $C_q:=2(s(A+2)+4)/\log2$, is the largest number of
predictors the conditions constrain jointly.

\Astate
\begin{localassumption}[Distribution-level setting]\label{asm:dist}
The integer $s\ge0$ and the constants $A\ge1$, $L\ge1$, $c_\psi>0$ and $\bar b>0$ are fixed.
\begin{asmitems}
  \item\label{asm:dist-model} \emph{Gaussian linear model for i.i.d.\ pairs.} $y=\alpha^*+x^\top\bstar+\varepsilon$ with
    $\varepsilon\sim\Normal(0,\sigs^2)$, $\sigs^2>0$, independent of $x$; $\alpha^*$, $\bstar$, $\sigs$ and the law of $x$ may
    depend on $n$.
  \item\label{asm:dist-sparse} \emph{Sparsity and dimension.} At most $s$ coordinates of $\bstar$ are nonzero, and
    $2\le\praw\le n^A$.
  \item\label{asm:dist-subg} \emph{Sub-Gaussian predictors on sparse directions.} Every predictor has positive variance,
    and $\E\exp(\lambda u^\top\zeta)\le\exp\big(L^2\lambda^2\Var(u^\top\zeta)/2\big)$ for every $\lambda\in\R$ and every
    $u\in\R^{\praw}$ with at most $q_n$ nonzero entries.
  \item\label{asm:dist-nec} \emph{No near-exact collinearity.} $\lmin(\Corr(x_T))\ge c_\psi/n$ for every set $T$ of at
    most $q_n$ predictors.
  \item\label{asm:dist-bs} \emph{Bounded signal.} $\abs{\beta^*_j}\sd(x_j)\le\bar b\,\sigs$ for every $j$.
\end{asmitems}
\end{localassumption}
\Awords{Items (a)--(c) are \as{GL}, \as{DIM}, \as{SP} and \as{SGR}, relaxed: the noise level and the coefficients may vary
with $n$, and no coefficient need be nonzero. Item (d) asks only that no small group of predictors be \emph{almost exactly}
collinear: its floor, $c_\psi/n$, tends to zero. Item (e) asks that no single predictor carry an unbounded signal relative
to the noise.}
\Arules{Duplicated columns and exact linear dependences (as in one-hot or compositional coding), for (d); a predictor
whose effect is unbounded in noise units, for (e).}
\Areas{Yes: every design we study is inside, and (d) is very mild; a user is inside after removing exact duplicates and
exact linear dependences. There is no lower bound on the coefficients, so weak effects are allowed. Without (d) the result
still holds with one extra factor $\sqrt{\log(\praw\vee n)}$ in the radius (supplement, Proposition~\ref{prop:dist}).}
\Asl{Items (a)--(c) and (e) compare as \as{GL}, \as{DIM}, \as{SP}, \as{SGR} and the signal bounds above. Item (d) is
\emph{mixed}: their Theorems~1--2 have no such condition, so there ours is \emph{stronger}; against their Condition~3
(p.~955) and (22) (supp.\ p.~13), if their sets are at least as large as ours, theirs implies ours, so ours is
\emph{weaker}. Our floor, $c_\psi/n$ on the correlation scale, is of the same order as the ridge $c_n\propto1/n$ they use
(p.~953).}

% ------------------------------------------------------------------------------------------------
\subsection{The assumptions of the reported-model result}\label{sec:a-rm}

The start law sees interactions within blocks of at most $16$ predictors only, so these conditions are about how the
signal looks given at most $15$ other predictors. With $W\sim\Normal_p(0,\Sig)$, $\Cov(W_j,W_l\mid W_U)$ is the partial
covariance given $(W_k)_{k\in U}$, and $\bar\beta_{U,l}:=\Sigma_{n,UU}^{-1}\Sigma_{n,Ul}$ is the population regression
coefficient of $W_l$ on $W_U$, padded with zeros to a vector indexed by a larger set when needed. A \emph{leakage set} is
a deterministic $Q$ with $\Sz\subseteq Q\subseteq\set{p}$ that contains every predictor able to imitate part of the
signal; its \emph{lattice} $\Lat(Q):=\{T:\Sz\subseteq T\subseteq Q\}$ is the set of models between the true model and $Q$.
In Example~\ref{ex:leak}, predictor $5$ leaks: it imitates part of predictor $1$'s signal.

\ahead{Visibility}{\lab{V}}{}
\Astate
\begin{assumption}{V}{Visibility: no true predictor masked by $\le15$ others}
For every $j\in\Sz$ and every $U\not\ni j$ with $\abs{U}\le15$, $\abs{\Cov(W_j,W^\top\tstar\mid W_U)}\ge a$, for a fixed
$a>0$.
\end{assumption}
\Awords{Each true predictor stays visibly related to the signal, by a fixed margin, whichever $15$ other predictors are
held fixed. A method that looks at blocks of $16$ must see every true predictor inside some block.}
\Arules{Masking: a true predictor whose signal some $15$ other predictors cancel. With two true predictors with
coefficients $1$ and $-1$ and correlation $\rho$, the signal of each, given nothing, is $1-\rho$, which vanishes as $\rho\to1$.}
\Areas{Yes in two broad cases (supplement, Remark~\ref{rem:vis}): (i) \emph{aligned signs}, for example independent
columns, autoregressive correlation, one factor with coefficients of one sign, or blocks with the true predictors in
different blocks; and (ii) \emph{generic coefficients}: for designs with finitely many local correlation patterns, every
coefficient vector outside one fixed null set of hyperplanes. It fails exactly on cancellations, such as adjacent AR(1)
true predictors at correlation $0.5$ with coefficients $1$ and $-2$. Inside: our independent and locally correlated
designs and every design of the strong-factor study. Unclear: S5 and FAC (mixed signs; a search found a near-cancellation
in S5). A design modelled on the riboflavin gene-expression data is inside when its four true predictors are spread apart
and not when they are clustered.}
\Asl{\emph{Stronger.} Ours implies their Condition~5 (p.~958: $\abs{\operatorname{cov}(x_t,y)}>0$ for every true
predictor) in every regime; theirs is ours with no conditioning and no margin.}

\ahead{Small leakage set with local regressions}{Assumption~\ref{asm:reported}}{}
Fix $b\in(0,\tfrac12)$ and $b_H>0$ with $15b_H+b<\tfrac12$.
\Astate
\begin{localassumption}[Polynomial leakage set with local regressions]\label{asm:reported}
\begin{asmitems}
  \item\label{asm:rr-leak} \emph{Leakage set.} There is a deterministic $Q\supseteq\Sz$ with $\abs{Q}\le n^b$ and, for
    large $n$, $\abs{Q}\le K$ and ($Q=\set{p}$ or $K\ge2\abs{Q}+32\sz+1$), such that: (i) \emph{no leak outside $Q$:}
    there is a deterministic $\varpi_n$ with $\varpi_n\sqrt{n/\log n}\to0$ such that
    $\abs{\Cov(W_l,W_j\mid W_U)}\le\varpi_n$ for all $l\notin Q$, $j\in\Sz\setminus U$ and $\abs{U}\le15$; and (ii)
    \emph{locality:} there is a deterministic $r_n\to0$, and every $l\notin Q$ has a deterministic
    $\Hl_l\supseteq\{l\}\cup\Sz$ with $\abs{\Hl_l}\le n^{b_H}$, such that for all $l\notin Q$, all $U\not\ni l$ with
    $\abs{U}\le15$ and $U_0:=U\cap\Hl_l$: $\abs{\Cov(W_l,W_v\mid W_{U_0})}\le r_n$ and
    $\abs{\Cov(W_j,W_v\mid W_{U_0})}\le r_n$ for all $v\in U\setminus\Hl_l$ and $j\in\Sz\setminus U_0$.
  \item\label{asm:rr-floor} \emph{Floor on the leakage set.} $\lmin(\Sigma_{n,Q\cup\{l\}})\ge\varphi^2$ for every
    $l\in\set{p}$, with $\varphi$ as in \as{EFS}.
  \item\label{asm:rr-local} \emph{Local regressions within the leakage set.} There are deterministic sequences
    $d_n=o(\log n)$ and $L_n$, and for every $l\in\set{p}$ a deterministic $N_l\subseteq Q\setminus\{l\}$ with
    $\abs{N_l}\le d_n$, such that $\norm{\bar\beta_{T,l}-\bar\beta_{T\cap N_l,l}}_1\le L_n$ for every $T\in\Lat(Q)$ with
    $l\notin T$, where $L_n\to0$ and $L_n^2\log(\abs{Q}+1)=o(\log n)$.
\end{asmitems}
\end{localassumption}
\Awords{(a) Only a small set $Q$ of predictors, at most $n^b$ with $b<\tfrac12$, can imitate part of the signal; every
predictor outside $Q$ is, given up to $15$ others, nearly uncorrelated with the true ones, and only its nearby predictors
matter. (b) The leakage set is not collinear with any predictor. (c) Within $Q$, the regression of a predictor on the others
uses only a few of them. In~(a), $32\sz=2\cdot16\sz$: with probability tending to one the start mode has at most
$2(\abs{Q}+16\sz)$ predictors, so the size condition leaves room to add a missing true predictor to it.}
\Arules{Pervasive factors with $p\gg\sqrt n$: with no conditioning every null predictor then has a fixed covariance with
every true one, so the leakage set is all of $\set{p}$, and $\abs{Q}\le n^b$ would force $p\le n^b$. Also leakage sets
collinear with another predictor, or whose regressions use many members.}
\Areas{Yes for block and autoregressive correlation, the usual models of local dependence (genes in pathways, markers along
a chromosome): $Q$ is the true predictors' blocks (S3, S4, S6, BLK), or a neighbourhood of size $C\log n$ around them (S2,
AR); for independent predictors $Q=\Sz$ (S1, IND). Items (b) and (c) then follow from a floor $\lmin(\Sig)\ge\varphi^2$ and
the locality of $\Sig$. No for pervasive factors (S5, the strong-factor study, FAC). The bound $b<\tfrac12$ comes from the
proof: at model size about $\sqrt n$ the penalty of one more predictor stops dominating the best of $p$ noise gains; whether
the result fails for larger $b$ is open.}
\Asl{\emph{None.} \SL{} have no hypothesis of this kind. The nearest text is their informal remark~(a), p.~958 (``if there
are $p_n$ predictors correlated with the true predictor $X_t$, $p_n\prec P_n$, then there are at most $cp_n$ predictors to
be selected''), which is not a condition of any of their theorems; their results on their sampler's output carry no design
condition of this kind (supp.\ pp.~9--10, 12--13, 16--17). The floor~(b) compares as \as{EFS}; the local regressions~(c)
have no counterpart (they serve the method's local search, which \SL{} do not analyse).}

% ------------------------------------------------------------------------------------------------
\subsection{The assumptions of certified screening}\label{sec:a-scr}

\ahead{Floor on sets of the budget size, and fixed signal bounds}{(Theorem~\ref{thm:screen})}{}
\Astate The design is fixed; \as{GL} holds; for a model-size budget $k_n<K$ chosen by the analyst, $\lmin(\Rhat_S)\ge\kappa$
for every $S$ with $\abs{S}\le k_n+\sz$; $\tmin\ge\tlow>0$ and $\norm{\tstar}_1\le M_1$; with $\kappa$, $\tlow$, $M_1$,
$\sigs^2$ fixed; and $(k_n+1)\log p+\log n\le c_{\mathrm{SB}}\,n$ with $c_{\mathrm{SB}}:=\kappa^2\tlow^2/\{200(M_1^2+\sigs^2)\}$.
\Awords{The observed predictors have an eigenvalue floor on sets of the budget size; the true effects are bounded below
and their total size is bounded; and the sample is large compared with the budget times $\log p$.}
\Arules{Collinearity within sets of $k_n+\sz$ predictors, signals that vary with $n$, and budgets too large for the
sample.}
\Areas{The floor and signal bounds are of the same kind as \as{EFS} and \as{FSt}, on sets of the budget size, which can be
much larger than a multiple of the sparsity. The budget inequality is a sufficient condition for a union bound and is
conservative: at the constants of the strong-factor designs it asks $n\gtrsim1{,}450$, while in those designs the certified
model already contained every true predictor in all $30$ fits at $n=100$ (Remark~\ref{rem:screen}). No visibility or
leakage condition is needed, so pervasive factors are allowed.}
\Asl{\emph{As} the eigenvalue floor and the signal bounds above, on other sets. \SL's sure-screening Theorem~3 (p.~956) is
the nearest conclusion; Section~\ref{sec:proved} compares the two.}

% ------------------------------------------------------------------------------------------------
\subsection{Conventions: what ``the method'' means}\label{sec:conventions}

Results 3--5 are about runs of the software. They use three conventions, which fix what ``the method'' means and
idealise the computer; they are not statistical assumptions (Appendix~\ref{app:conventions} gives them in full).
\begin{itemize}[leftmargin=1.5em,itemsep=1pt]
  \item \phantomsection\label{asm:C1}\emph{Convention \as{C1} (the method):} exact arithmetic, the default configuration,
    valid inputs and seed. A run that meets it is \emph{admissible}; a sample size $n$ is admissible if some data set of
    that size is valid. Statements hold for every valid seed and every outcome of the calibration.
  \item \phantomsection\label{asm:C2p}\emph{Convention \as{C2p} (delivery):} an idealised
    machine that never overflows or fails, ideal random numbers, an allocation of at least the memory and time that the
    method's own polynomial requirement formulas require, and the code's two fixed integer limits on the move caps lifted
    (neither can be reached on any input R can represent).
  \item \phantomsection\label{asm:C3}\emph{Convention \as{C3} (sampler):} exact arithmetic, ideal random numbers
    independent of the data and of the calibration, an exact start sampler, and Markov moves that do what their definitions
    say.
\end{itemize}
The results also use the event $\Acal$ that the method uses the sampler and its calibration succeeds; every completed run
lies in $\Acal$. \emph{Reasonable?} These are the usual idealisations of any theorem about an algorithm. \emph{vs \SL:}
exact arithmetic and ideal random numbers are presupposed by their statements about their own sampler too (supp.,
pp.~9--10, 12--13).
